\documentclass[11pt, letterpaper]{article}


\usepackage[utf8]{inputenc}
\usepackage[spanish, es-tabla]{babel}
\usepackage[margin=2.3cm, top=2.8cm, bottom=2.8cm]{geometry}
\usepackage{helvet} 
\renewcommand{\familydefault}{\sfdefault}
\usepackage{microtype}
\usepackage{xcolor}
\usepackage{graphicx}
\usepackage{booktabs}
\usepackage{tabularx}
\usepackage{fancyhdr}
\usepackage{lastpage}
\usepackage{hyperref}
\usepackage{fontawesome5} % Para incluir el ícono de GitHub


\definecolor{usmblue}{RGB}{0, 75, 141}
\definecolor{usmred}{RGB}{200, 16, 46}
\definecolor{usmgray}{RGB}{100, 100, 100}
\definecolor{lightgray}{RGB}{245, 245, 245}


\pagestyle{fancy}
\fancyhf{}
\renewcommand{\headrulewidth}{0.4pt}
\renewcommand{\footrulewidth}{0.4pt}


\fancyhead[L]{\small \color{usmgray} Universidad Técnica Federico Santa María\\ \textbf{ICN-292 — Sistemas de Información para la Gestión}}
\fancyhead[R]{\small \color{usmgray} Departamento de Industrias\\ Segundo Semestre 2026}


\fancyfoot[L]{\small \color{usmgray} Laboratorio 4 · ICN-292 · Informe Oficial}
\fancyfoot[R]{\small \color{usmgray} Página \thepage\ de \pageref{LastPage}}

\hypersetup{
    colorlinks=true,
    linkcolor=usmblue,
    urlcolor=usmblue,
    pdftitle={Laboratorio 1 - ICN292}
}

\begin{document}


\thispagestyle{empty}


\begin{raggedright}
    {\bfseries\large UNIVERSIDAD TECNICA FEDERICO SANTA MARIA}\\[0.1cm]
    {\large Departamento de Industrias}\\[0.3cm]
    {\color{usmblue} \rule{\textwidth}{1.5pt}}\\[0.4cm]
    {\bfseries\color{usmblue}\huge INFORME DE LABORATORIO 4}\\[0.2cm]
    {\Large Entregable — SQL (1/2)
}\\[0.1cm]
    {\large\textbf{ICN-292 — Sistemas de Información para la Gestión}}\\[0.1cm]
    {\large Segundo Semestre 2026}
\end{raggedright}

\vspace{0.3cm}


\begin{center}
\colorbox{lightgray}{
\begin{minipage}{0.95\textwidth}
\vspace{0.2cm}
\begin{tabular}{ll}
    \textbf{Estudiante:} & Martin Ariel Aspee Arroyo \\
   
    \textbf{Fecha de Entrega:} & Sábado 7 de Nopviembre de 2026, 11:00 hrs \\
    \textbf{Campus / Sede:} & Vitacura \\
    \textbf{Profesores:} & José Miguel González / Paul José Luis Sáez Tamayo
\end{tabular}
\vspace{0.2cm}
\end{minipage}
}
\end{center}


\vspace{0.3cm}
\begin{center}
    \href{https://github.com/Diaspuu/ICN292-Lab4-Aspee-Martin}{\Large \faGithub} \\
    \vspace{0.1cm}
    \url{https://github.com/Diaspuu/ICN292-Lab4-Aspee-Martin}
\end{center}
\vspace{0.3cm}
\section{diagrama E-R}
\begin{figure}[h]
    \centering
    \includegraphics[width=0.6\linewidth]{diagrama e-r.png}
    \caption{Diagrama de entidad-relación al hacer ingenieria inversa en SQL workbench}
    \label{fig:placeholder}
\end{figure}
\section{Consultas SQL a desarrollar}

\subsection{consulta 1}

\begin{figure} [h]
    \centering
    \includegraphics[width=0.5\linewidth]{image.png}
    \caption{Código SQL y resultado entregado}
    \label{fig:placeholder}
\end{figure}
La comparacion devuelve todos los paises que cumplen la desigualdad dejando de lado al pais 'null' puesto que la entrada 'null' no se considera igual o mayor/menor a otro valor.

\subsection{consulta 2}

\begin{figure} [h]
    \centering
    \includegraphics[width=0.5\linewidth]{consulta 2.png}
    \caption{Código SQL y resultado entregado}
    \label{fig:placeholder}
\end{figure}
\newpage
\subsection{consulta 3}

\begin{figure}[h]
    \centering
    \includegraphics[width=0.5\linewidth]{consulta 3.png}
    \caption{Código SQL y resultado entregado}
    \label{fig:placeholder}
\end{figure}

\subsection{consulta 4}
\begin{figure}[h]
    \centering
    \includegraphics[width=0.5\linewidth]{consulta 4.png}
    \caption{Código SQL y resultado entregado}
    \label{fig:placeholder}
\end{figure}
\newpage
\subsection{consulta 5}
\begin{figure}[h]
    \centering
    \includegraphics[width=0.5\linewidth]{consulta 5.png}
    \caption{Código SQL y resultado entregado}
    \label{fig:placeholder}
\end{figure}

\subsection{consulta 6}
\begin{figure}[h]
    \centering
    \includegraphics[width=0.4\linewidth]{consulta 6.png}
    \caption{Código SQL y resultado entregado}
    \label{fig:placeholder}
\end{figure}
\newpage
\subsection{consulta 7}
\begin{figure}[h]
    \centering
    \includegraphics[width=0.5\linewidth]{consulta 7.png}
    \caption{Código SQL y lista entregado al agregar el libro nuevo}
    \label{fig:placeholder}
\end{figure}
\begin{figure}[h]
    \centering
    \includegraphics[width=0.5\linewidth]{consulta 7_2.png}
    \caption{Código SQL y lista entregado al eliminar el libro nuevo}
    \label{fig:placeholder}
\end{figure}
\newpage
\subsection{consulta 8}
\begin{figure}[h]
    \centering
    \includegraphics[width=0.5\linewidth]{consulta 8.png}
    \caption{Código SQL y resultado entregado}
    \label{fig:placeholder}
\end{figure}

\subsection{consulta 9}
\begin{figure}[h]
    \centering
    \includegraphics[width=0.5\linewidth]{consulta 9.png}
    \caption{Código SQL y resultado entregado}
    \label{fig:placeholder}
\end{figure}
\newpage
\subsection{consulta 10}
\begin{figure}[h]
    \centering
    \includegraphics[width=0.5\linewidth]{consulta 10 1.png}
    \caption{Código SQL y resultado entregado}
    \label{fig:placeholder}
\end{figure}
\begin{figure}[h]
    \centering
    \includegraphics[width=0.5\linewidth]{consulta 10 2.png}
    \caption{Código SQL y resultado entregado}
    \label{fig:placeholder}
\end{figure}
\end{document}